\chapter{Alexander Grothendieck --- Fields Medal 1966}\label{ch1966alexandergrothendieck:alexander-grothendieck-fields-medal-1966-a-comprehensive-report}

\begin{center}\rule{0.5\linewidth}{0.5pt}\end{center}

\textbf{Prepared:} July 2026 \textbf{Subject:} Alexander Grothendieck
(b. March 28, 1928, Berlin, Germany) \textbf{Position:} Professor, IHÉS;
University of Montpellier \textbf{Doctoral Advisor:} Laurent Schwartz
(Université Henri Poincaré, Nancy) \textbf{Fields Medal Citation:}
\emph{``Built on work of Weil and Zariski and effected fundamental
advances in algebraic geometry. He introduced the idea of K-theory (the
Grothendieck groups and rings). Revolutionized homological algebra in
his well-known `Tohoku paper'.''}

\begin{center}\rule{0.5\linewidth}{0.5pt}\end{center}

\section{1. Introduction}\label{ch1966alexandergrothendieck:introduction}

Alexander Grothendieck was awarded the Fields Medal in 1966 for his
significant contributions to algebraic geometry, homological algebra,
and functional analysis. He is a prominent mathematician of the 20th century, renowned for his notable
vision and his ability to rebuild entire fields from their foundations.

Grothendieck's work transformed algebraic geometry into a unified,
powerful subject through the theory of \textbf{schemes}, which provides
a common framework for algebraic varieties over any field or ring. He
developed \textbf{étale cohomology}, a Weil cohomology theory that
allowed the proof of the Weil conjectures (completed by Deligne). He
created the theory of \textbf{motives}, a universal cohomology theory
for algebraic varieties. He revolutionized homological algebra with the
\textbf{Tohoku paper}, which developed abelian categories and derived
functors. He introduced \textbf{K-theory} in the context of the
Grothendieck--Riemann--Roch theorem.

Grothendieck's major contributions include: - \textbf{Schemes}: The
foundation of modern algebraic geometry. - \textbf{Étale cohomology and
the Weil conjectures}. - \textbf{Motives}: The universal cohomology
theory. - \textbf{The Tohoku paper}: Abelian categories and derived
functors. - \textbf{Grothendieck--Riemann--Roch theorem}. -
\textbf{Topos theory}: A generalization of topological spaces. -
\textbf{Dessins d'enfants}: Galois action on combinatorial objects. -
\textbf{Grothendieck's Galois theory}: A categorical generalization of
classical Galois theory.

\begin{center}\rule{0.5\linewidth}{0.5pt}\end{center}

\section{2. Biographical Background}\label{ch1966alexandergrothendieck:biographical-background}

\subsection{2.1 Early Life and
Education}\label{ch1966alexandergrothendieck:early-life-and-education}

Alexander Grothendieck was born on March 28, 1928, in Berlin to
anarchist parents. His father, Alexander ``Sascha'' Schapiro, was a
Russian Jew, and his mother, Johanna ``Hanka'' Grothendieck, was German.
During the Nazi era, his father was interned and eventually killed at
Auschwitz. Alexander lived in hiding in France during World War II.

After the war, he studied mathematics at the University of Montpellier,
then at the École Normale Supérieure in Paris, and finally at the
University of Nancy, where he completed his doctorate under Laurent
Schwartz.

\subsection{2.2 Academic Career}\label{ch1966alexandergrothendieck:academic-career}

\begin{itemize}
\tightlist
\item
 \textbf{1945--1948:} Studies at the University of Montpellier.
\item
 \textbf{1948--1949:} Studies at the École Normale Supérieure.
\item
 \textbf{1949--1953:} Studies and Ph.D.~at Nancy. Thesis: ``Produits
 tensoriels topologiques et espaces nucléaires.''
\item
 \textbf{1953--1955:} CNRS researcher, Paris.
\item
 \textbf{1955--1957:} Visiting professor, University of São Paulo;
 University of Kansas.
\item
 \textbf{1957--1960:} Professor, Institut Henri Poincaré.
\item
 \textbf{1959--1970:} Professor, IHÉS (Institut des Hautes Études
 Scientifiques).
\item
 \textbf{1966:} Fields Medal at the ICM in Moscow.
\item
 \textbf{1970--1973:} Professor, Collège de France.
\item
 \textbf{1973--1988:} Professor, University of Montpellier.
\item
 \textbf{1991:} Retired to the village of Lasserre in the Pyrenees.
\item
 \textbf{2014:} Died at age 86 in Saint-Girons, France.
\end{itemize}

\subsection{2.3 Doctoral Work}\label{ch1966alexandergrothendieck:doctoral-work}

Grothendieck's early work was in functional analysis. His thesis on
topological tensor products and nuclear spaces (1953) was a major
contribution to the field. He introduced the concept of \textbf{nuclear
spaces}, which became fundamental in the theory of distributions and in
the structure theory of topological vector spaces.

\begin{center}\rule{0.5\linewidth}{0.5pt}\end{center}

\section{3. The Tohoku Paper}\label{ch1966alexandergrothendieck:the-tohoku-paper}

\subsection{3.1 Background}\label{ch1966alexandergrothendieck:background}

In the 1950s, homological algebra was being developed by Cartan,
Eilenberg, and Mac Lane. The theory was powerful but somewhat ad hoc,
dealing separately with modules over rings, sheaves on topological
spaces, and other categories.

\subsection{3.2 Grothendieck's
Contribution}\label{ch1966alexandergrothendieck:grothendiecks-contribution}

In his 1957 paper ``Sur quelques points d'algèbre homologique''
(published in the Tohoku Mathematical Journal), Grothendieck established
the modern framework of homological algebra. He introduced: -
\textbf{Abelian categories}: categories with kernels, cokernels, and
exactness properties. - \textbf{Grothendieck axioms}: AB1--AB6 for
abelian categories. - \textbf{Injective envelopes}: existence of minimal
injective extensions in any Grothendieck category (AB5). -
\textbf{Derived functors}: defined using injective (or projective)
resolutions in any abelian category with enough injectives.

The key axiom AB5 states that a filtered colimit of exact sequences is
exact, which holds for categories of modules and for sheaves on a space.
Grothendieck proved that any AB5 category with a generator has enough
injectives, providing the theoretical foundation for sheaf cohomology.

\subsection{3.3 Impact}\label{ch1966alexandergrothendieck:impact}

The Tohoku paper is a widely cited paper in mathematics. It
provided the categorical framework that underlies modern algebraic
geometry, algebraic topology, and representation theory. Every
mathematician who uses sheaf cohomology, derived categories, or
homological algebra is building on Grothendieck's foundations.

\begin{center}\rule{0.5\linewidth}{0.5pt}\end{center}

\section{4. Schemes}\label{ch1966alexandergrothendieck:schemes}

\subsection{4.1 The Concept}\label{ch1966alexandergrothendieck:the-concept}

A \textbf{scheme} is a locally ringed space \((X, \mathcal{O}_X)\)
modeled on the spectrum of a ring. The \textbf{prime spectrum}
\(\text{Spec}(R)\) of a commutative ring \(R\) is the set of prime
ideals of \(R\), with the Zariski topology and a sheaf of rings
\(\mathcal{O}_{\text{Spec}(R)}\).

A scheme \(X\) is a topological space with a sheaf of commutative rings
\(\mathcal{O}_X\) (the structure sheaf) such that every point has an
open neighborhood \(U\) with
\((U, \mathcal{O}_X|_U) \cong \text{Spec}(R)\) for some ring \(R\).

\subsection{4.2 The Grothendieck
Revolution}\label{ch1966alexandergrothendieck:the-grothendieck-revolution}

Grothendieck introduced schemes in the monumental \textbf{Éléments de
Géométrie Algébrique} (EGA), written with Jean Dieudonné. The scheme
perspective unifies: - \textbf{Classical algebraic varieties}: over
algebraically closed fields. - \textbf{Arithmetic geometry}: varieties
over \(\mathbb{Z}\), \(\mathbb{Q}\), and number fields. -
\textbf{Deformation theory}: infinitesimal neighborhoods using dual
numbers \(k[\epsilon]/(\epsilon^2)\). - \textbf{Moduli spaces}: families
of algebraic objects.

\subsection{4.3 Functorial Point of
View}\label{ch1966alexandergrothendieck:functorial-point-of-view}

Grothendieck emphasized that a scheme \(X\) is best understood by its
\textbf{functor of points}: for any scheme \(S\), the set
\(X(S) = \text{Hom}_{\text{Sch}}(S, X)\) parametrizes families of points
of \(X\) over \(S\). This functorial perspective is essential for moduli
theory and the study of algebraic stacks.

\subsection{4.4 Fibered Products}\label{ch1966alexandergrothendieck:fibered-products}

Schemes admit \textbf{fibered products} \(X \times_S Y\), which
generalize the product of varieties. This allows the systematic study of
families: a morphism \(X \to S\) is a family of schemes parametrized by
\(S\). Properties of morphisms (smooth, proper, flat, étale) become
central to the theory.

\begin{center}\rule{0.5\linewidth}{0.5pt}\end{center}

\section{5. Étale Cohomology and the Weil
Conjectures}\label{ch1966alexandergrothendieck:uxe9tale-cohomology-and-the-weil-conjectures}

\subsection{5.1 Motivation}\label{ch1966alexandergrothendieck:motivation}

The \textbf{Weil conjectures} (1949) predicted deep properties of zeta
functions of algebraic varieties over finite fields. They required a
cohomology theory for varieties over fields of characteristic \(p\) that
would behave like singular cohomology for complex manifolds. Ordinary
sheaf cohomology in the Zariski topology is inadequate because the
Zariski topology is too coarse.

\subsection{5.2 The Étale Topology}\label{ch1966alexandergrothendieck:the-uxe9tale-topology}

Grothendieck's solution was the \textbf{étale topology}, a refinement of
the Zariski topology where coverings are given by étale morphisms
(smooth maps of relative dimension zero, analogous to local isomorphisms
in complex analysis).

The \textbf{étale cohomology groups}
\(H^i_{\text{ét}}(X, \mathbb{Q}_\ell)\) are defined using sheaves on the
étale site. They satisfy: - \textbf{Finite fields}:
\(H^i_{\text{ét}}(X, \mathbb{Q}_\ell)\) is a finite-dimensional
\(\mathbb{Q}_\ell\)-vector space. - \textbf{Poincaré duality}: for
smooth proper varieties. - \textbf{Lefschetz fixed point formula}: for a
morphism \(f: X \to X\). - \textbf{Comparison theorem}: for complex
varieties, \(H^i_{\text{ét}}(X, \mathbb{Q}_\ell) \cong
H^i(X(\mathbb{C}), \mathbb{Q}) \otimes \mathbb{Q}_\ell\).

\subsection{5.3 The Proof of the Weil
Conjectures}\label{ch1966alexandergrothendieck:the-proof-of-the-weil-conjectures}

Grothendieck proved the first three Weil conjectures using étale
cohomology: 1. \textbf{Rationality}: the zeta function is a rational
function. 2. \textbf{Functional equation}: symmetry under \(q^{-s}\). 3.
\textbf{Riemann hypothesis} (partial): Deligne completed the proof of
the Riemann hypothesis part in 1974.

\subsection{5.4 The Standard
Conjectures}\label{ch1966alexandergrothendieck:the-standard-conjectures}

Grothendieck proposed the \textbf{Standard Conjectures} about the
behavior of algebraic cycles under the Lefschetz and Künneth operators.
These conjectures remain largely open, but they are fundamental to the
theory of motives.

\begin{center}\rule{0.5\linewidth}{0.5pt}\end{center}

\section{6. The Grothendieck--Riemann--Roch
Theorem}\label{ch1966alexandergrothendieck:the-grothendieckriemannroch-theorem}

\subsection{6.1 Background}\label{ch1966alexandergrothendieck:background-1}

The \textbf{Hirzebruch--Riemann--Roch theorem} expresses the Euler
characteristic of a holomorphic vector bundle on a projective manifold
in terms of Chern classes. Grothendieck realized this was a statement
about the behavior of the Chern character under pushforward maps.

\subsection{6.2 The Theorem}\label{ch1966alexandergrothendieck:the-theorem}

For a proper morphism \(f: X \to Y\) of smooth quasi-projective
varieties, the Grothendieck--Riemann--Roch theorem states:

\[\text{ch}(f_!(E)) \cdot \text{td}(T_Y) = f_*(\text{ch}(E) \cdot \text{td}(T_X))\]

Here \(f_!\) is the direct image in K-theory, \(f_*\) is the direct
image in homology, \(\text{ch}\) is the Chern character, and
\(\text{td}\) is the Todd class.

The proof uses the concept of \textbf{K-theory}, which Grothendieck
introduced for this purpose. The \textbf{Grothendieck group} \(K(X)\) of
a scheme \(X\) is the free abelian group on isomorphism classes of
vector bundles, modulo exact sequences.

\subsection{6.3 Impact}\label{ch1966alexandergrothendieck:impact-1}

The Grothendieck--Riemann--Roch theorem is fundamental in algebraic
geometry. It implies the original Hirzebruch--Riemann--Roch theorem, and
it has applications to: - The computation of Chern numbers of varieties.
- Moduli theory: the geometry of the moduli space of curves and vector
bundles. - The structure of algebraic cycles and the theory of motives.

\begin{center}\rule{0.5\linewidth}{0.5pt}\end{center}

\section{7. Motives}\label{ch1966alexandergrothendieck:motives}

\subsection{7.1 The Concept}\label{ch1966alexandergrothendieck:the-concept-1}

Grothendieck envisioned a universal cohomology theory for algebraic
varieties, the theory of \textbf{motives}. Intuitively, a motive is an
abstract algebraic object that captures the ``essential geometric
content'' of a variety, independent of which particular cohomology
theory you use to study it. Just as a physical object can be described
by its color, shape, or temperature, a variety can be studied through
different cohomological ``descriptions''; the motive is the underlying
object itself.

More precisely, a motive \(M(X)\) is a
universal object from which all cohomology theories (singular, étale, de
Rham, Hodge, crystalline) can be obtained by realization functors.

The category of (pure) motives is built from: - Objects: smooth
projective varieties \(X\). - Morphisms: correspondences modulo rational
equivalence. - The \textbf{Tate motive} \(\mathbb{Q}(1)\): the Lefschetz motive, providing the twisting.

\subsection{7.2 The Motivic Galois
Group}\label{ch1966alexandergrothendieck:the-motivic-galois-group}

Grothendieck conjectured the existence of a \textbf{motivic Galois
group} that acts on all cohomology realizations of motives. This would
explain the relationship between different cohomology theories and their
symmetry groups.

\subsection{7.3 Impact}\label{ch1966alexandergrothendieck:impact-2}

The theory of motives remains an active area of algebraic
geometry. Voevodsky (Fields Medal 2002) developed \textbf{triangulated
categories of motives} and proved the Milnor conjectures using them. The
theory connects to the Langlands program, the Birch--Swinnerton-Dyer
conjecture, and the theory of automorphic forms.

\begin{center}\rule{0.5\linewidth}{0.5pt}\end{center}

\section{8. Topos Theory}\label{ch1966alexandergrothendieck:topos-theory}

\subsection{8.1 Definition}\label{ch1966alexandergrothendieck:definition}

A \textbf{topos} is a category that behaves like the category of sheaves
on a topological space. Grothendieck defined a \textbf{topos} as the
category of sheaves on a site (a category with a Grothendieck topology).
He developed this theory in \emph{Séminaire de Géométrie Algébrique}
(SGA 4).

\subsection{8.2 Grothendieck Topologies}\label{ch1966alexandergrothendieck:grothendieck-topologies}

A \textbf{Grothendieck topology} on a category \(\mathcal{C}\)
specifies, for each object \(U\), a collection of families
\(\{U_i \to U\}\) (coverings) satisfying: - Isomorphisms are coverings.
- Coverings pull back. - Coverings compose.

The category \(\text{Sh}(\mathcal{C})\) of sheaves on a site is a topos.

\subsection{8.3 Properties of Topoi}\label{ch1966alexandergrothendieck:properties-of-topoi}

Topoi have notable properties: - They are Cartesian closed. - They
have a subobject classifier \(\Omega\) (the truth value object). - They
satisfy the axioms of intuitionistic set theory. - Geometric morphisms
\(f: \mathcal{E} \to \mathcal{F}\) model continuous maps.

\subsection{8.4 Impact}\label{ch1966alexandergrothendieck:impact-3}

Topos theory unifies topology, algebraic geometry, and logic.
\textbf{Lawvere and Tierney} developed \textbf{elementary topoi},
leading to applications in set theory (the independence of CH via
Boolean topoi) and computer science (type theory).

\begin{center}\rule{0.5\linewidth}{0.5pt}\end{center}

\section{9. Descent Theory}\label{ch1966alexandergrothendieck:descent-theory}

\subsection{9.1 The Problem}\label{ch1966alexandergrothendieck:the-problem}

\textbf{Descent theory} asks: given a morphism \(f: S' \to S\) and an
object (sheaf, vector bundle, etc.) on \(S'\), when does it descend to
\(S\)? This is essential for patching local data into global objects.

\subsection{9.2 Grothendieck's
Contribution}\label{ch1966alexandergrothendieck:grothendiecks-contribution-1}

Grothendieck gave a complete theory of descent using the \textbf{Čech
nerve} of a covering. For an fpqc (faithfully flat quasi-compact)
morphism \(f: S' \to S\), the category of quasi-coherent sheaves on
\(S\) is equivalent to the category of descent data on \(S'\). This is
the \textbf{Grothendieck fpqc descent theorem}.

\subsection{9.3 Applications}\label{ch1966alexandergrothendieck:applications}

Descent theory is fundamental for: - \textbf{Algebraic stacks}: objects
that satisfy descent for étale or fppf coverings. - \textbf{Galois
descent}: the theory of forms and the Galois cohomology of algebraic
groups. - \textbf{The theory of bundles}: principal bundles and their
classification.

\begin{center}\rule{0.5\linewidth}{0.5pt}\end{center}

\section{10. The SGA Séminaires}\label{ch1966alexandergrothendieck:the-sga-suxe9minaires}

\subsection{10.1 The IHÉS Séminaires}\label{ch1966alexandergrothendieck:the-ihuxe9s-suxe9minaires}

At IHÉS, Grothendieck led the notable Séminaire de Géométrie
Algébrique (SGA), which produced volumes that are the foundational texts
of modern algebraic geometry: - \textbf{SGA 1}: Fundamental group,
covering spaces, dessins d'enfants. - \textbf{SGA 2}: Cohomology of
schemes, local cohomology. - \textbf{SGA 3}: Group schemes. -
\textbf{SGA 4}: Topos theory, étale cohomology. - \textbf{SGA 4½}:
Deligne's completion of the Weil conjectures. - \textbf{SGA 5}: The
Lefschetz formula. - \textbf{SGA 6}: K-theory and the Riemann--Roch
theorem. - \textbf{SGA 7}: Monodromy, Hodge theory.

\subsection{10.2 EGA}\label{ch1966alexandergrothendieck:ega}

The \textbf{Éléments de Géométrie Algébrique} (EGA) was intended as the
definitive treatise on scheme theory. The published volumes (with
Dieudonné) cover: - EGA I: The language of schemes. - EGA II: Étale and
smooth morphisms. - EGA III: Cohomology of schemes. - EGA IV: Flat and
proper morphisms.

EGA remains the standard reference for scheme theory.

\begin{center}\rule{0.5\linewidth}{0.5pt}\end{center}

\section{11. Dessins d'Enfants}\label{ch1966alexandergrothendieck:dessins-denfants}

\subsection{11.1 Definition}\label{ch1966alexandergrothendieck:definition-1}

A \textbf{dessin d'enfant} (child's drawing) is a finite graph drawn on
a compact orientable surface that cuts it into simply connected faces.
Equivalently, it is a bipartite graph embedded in a compact Riemann
surface.

\subsection{11.2 Belyi's Theorem and Galois
Action}\label{ch1966alexandergrothendieck:belyis-theorem-and-galois-action}

\textbf{Belyi's theorem} states that a Riemann surface \(X\) can be
defined over \(\overline{\mathbb{Q}}\) if and only if there exists a
non-constant morphism \(X \to \mathbb{P}^1\) ramified over at most three
points.

Such covers correspond bijectively to dessins d'enfants. Grothendieck
realized that this gives a concrete description of the action of the
absolute Galois group \(\text{Gal}(\overline{\mathbb{Q}}/\mathbb{Q})\)
on combinatorial objects.

\subsection{11.3 Impact}\label{ch1966alexandergrothendieck:impact-4}

Dessins d'enfants are a concrete window into the
structure of the absolute Galois group. They connect algebraic geometry,
combinatorics, and number theory in a deep way.

\begin{center}\rule{0.5\linewidth}{0.5pt}\end{center}

\section{12. Grothendieck's Galois
Theory}\label{ch1966alexandergrothendieck:grothendiecks-galois-theory}

\subsection{12.1 Categorical Galois
Theory}\label{ch1966alexandergrothendieck:categorical-galois-theory}

Grothendieck reformulated classical Galois theory in categorical terms:
For a field \(K\) with separable closure \(K^{\text{sep}}\), the
\textbf{Galois category} of finite étale \(K\)-algebras is equivalent to
the category of finite \(\text{Gal}(K^{\text{sep}}/K)\)-sets.

This reformulation generalizes to: - The fundamental group of a scheme.
- The étale fundamental group \(\pi_1^{\text{ét}}(X)\). - The theory of
covering spaces.

\subsection{12.2 The Étale Fundamental
Group}\label{ch1966alexandergrothendieck:the-uxe9tale-fundamental-group}

For a connected scheme \(X\) with a geometric point \(\bar x\), the
\textbf{étale fundamental group} \(\pi_1^{\text{ét}}(X, \bar x)\)
classifies finite étale coverings of \(X\). For complex varieties, it is
the profinite completion of the topological fundamental group. This
connects algebraic geometry to group theory and is essential for the
study of Galois representations.

\begin{center}\rule{0.5\linewidth}{0.5pt}\end{center}

\section{13. Nuclear Spaces and Functional
Analysis}\label{ch1966alexandergrothendieck:nuclear-spaces-and-functional-analysis}

\subsection{13.1 Nuclear Spaces}\label{ch1966alexandergrothendieck:nuclear-spaces}

Grothendieck's early work in functional analysis introduced
\textbf{nuclear spaces}: locally convex topological vector spaces where
every continuous linear map to a Banach space is a nuclear operator.
Examples: Schwartz space \(\mathcal{S}(\mathbb{R}^n)\), the space of
distributions \(\mathcal{D}'(\mathbb{R}^n)\), the space of holomorphic
functions \(\mathcal{O}(\Omega)\).

\subsection{13.2 The Kernel Theorem}\label{ch1966alexandergrothendieck:the-kernel-theorem}

The \textbf{Grothendieck kernel theorem} states that every continuous
linear map from a nuclear space \(E\) to the dual of another nuclear
space \(F'\) is given by a distributional kernel. This unifies the
theory of integral operators and is fundamental in the theory of PDEs.

\subsection{13.3 Tensor Products}\label{ch1966alexandergrothendieck:tensor-products}

Grothendieck developed the theory of topological tensor products,
introducing the projective and injective tensor products of locally
convex spaces. His work on tensor products laid the foundation for the
modern theory of operator spaces and quantum information theory.

\begin{center}\rule{0.5\linewidth}{0.5pt}\end{center}

\section{14. Grothendieck's Later Work and
Retirement}\label{ch1966alexandergrothendieck:grothendiecks-later-work-and-retirement}

\subsection{14.1 Récoltes et Semailles}\label{ch1966alexandergrothendieck:ruxe9coltes-et-semailles}

After leaving IHÉS in 1970, Grothendieck wrote the massive manuscript
\emph{Récoltes et Semailles} (Harvests and Sowings), a combination of
mathematical reflection, autobiography, and philosophical meditation. In
it, he reflects on his mathematical work, his relationships with
colleagues, and his vision for mathematics.

\subsection{14.2 La Clef des Songes}\label{ch1966alexandergrothendieck:la-clef-des-songes}

He also wrote \emph{La Clef des Songes} (The Key to Dreams), a work on
the interpretation of dreams and his personal spiritual journey. He
became increasingly isolated from the mathematical community, retiring
to the village of Lasserre in the Pyrenees in 1991.

\subsection{14.3 Mathematical Legacy}\label{ch1966alexandergrothendieck:mathematical-legacy}

Despite his withdrawal, Grothendieck's influence on mathematics is
substantial. His concepts --- schemes, étale cohomology, motives, topos
theory, K-theory, abelian categories --- are now part of the standard
language of algebraic geometry and beyond.

\begin{center}\rule{0.5\linewidth}{0.5pt}\end{center}

\section{15. Awards and Recognition}\label{ch1966alexandergrothendieck:awards-and-recognition}

\textbf{1966:} - Fields Medal (International Mathematical Union).

\textbf{1977:} - Émile Picard Medal (French Academy of Sciences).

\textbf{1988:} - Crafoord Prize (jointly with Pierre Deligne). (He
declined the prize.)

\textbf{Elected to:} - French Academy of Sciences (rejected membership).
- American Academy of Arts and Sciences (Foreign Honorary). - National
Academy of Sciences (Foreign Associate).

\begin{center}\rule{0.5\linewidth}{0.5pt}\end{center}

\section{16. Selected Major
Publications}\label{ch1966alexandergrothendieck:selected-major-publications}

\begin{enumerate}
\def\labelenumi{\arabic{enumi}.}
\item
 \textbf{``Sur quelques points d'algèbre homologique.''} \emph{Tohoku
 Mathematical Journal} 9, 119--221, 1957. --- \emph{The Tohoku paper,
 founding modern homological algebra.}
\item
 \textbf{``Éléments de géométrie algébrique I--IV.''} (with J.
 Dieudonné) \emph{Publications Mathématiques de l'IHÉS} 4, 8, 11, 17,
 20, 24, 28, 32, 1960--1967. --- \emph{The foundations of scheme
 theory.}
\item
 \textbf{``Séminaire de géométrie algébrique (SGA) 1--7.''}
 \emph{Lecture Notes in Mathematics} 224, 269, 270, 288, 305, 340, 569,
 Springer, 1971--1977. --- \emph{Foundational texts on étale
 cohomology, topos theory, K-theory, and the fundamental group.}
\item
 \textbf{``Produits tensoriels topologiques et espaces nucléaires.''}
 \emph{Memoirs of the AMS} 16, 1955. --- \emph{Thesis on nuclear spaces
 and topological tensor products.}
\item
 \textbf{``La théorie des classes de Chern.''} \emph{Bulletin de la
 Société Mathématique de France} 86, 137--154, 1958. --- \emph{The
 Grothendieck--Riemann--Roch theorem.}
\item
 \textbf{``Esquisse d'un programme.''}
\end{enumerate}

\begin{enumerate}
\def\labelenumi{(\arabic{enumi})}
\tightlist
\item
 Published in \emph{Geometric Galois Actions I}, Cambridge, 1997. ---
 \emph{Outlines of his program on dessins d'enfants, anabelian
 geometry, and motives.}
\end{enumerate}

\begin{enumerate}
\def\labelenumi{\arabic{enumi}.}
\item
 \textbf{``Récoltes et semailles.''} Manuscript, 1986. ---
 \emph{Autobiographical and mathematical reflections.}
\item
 \textbf{``La longue marche ininterrompue de la géométrie
 algébrique.''} Manuscript, 1987. --- \emph{Philosophical reflections
 on algebraic geometry.}
\item
 \textbf{``Technique de descente et théorèmes d'existence en géométrie
 algébrique I--VI.''} \emph{Séminaire Bourbaki}, 1959--1962. ---
 \emph{Descent theory and existence theorems.}
\item
 \textbf{``Groupes de Barsotti--Tate et cristaux de Dieudonné.''}
 (1986) --- \emph{The theory of \(p\)-divisible groups and crystals.}
\end{enumerate}

\begin{center}\rule{0.5\linewidth}{0.5pt}\end{center}

\section{17. References and Further
Reading}\label{ch1966alexandergrothendieck:references-and-further-reading}

\subsection{Official Sources}\label{ch1966alexandergrothendieck:official-sources}

\begin{itemize}
\tightlist
\item
 IMU Fields Medal 1966:
 https://www.mathunion.org/imu-awards/fields-medal/fields-medals-1966
\item
 Crafoord Prize: https://www.crafoordprize.se/
\end{itemize}

\subsection{Key Technical References}\label{ch1966alexandergrothendieck:key-technical-references}

\begin{itemize}
\tightlist
\item
 Grothendieck, A. \emph{EGA I--IV}. IHÉS, 1960--1967.
\item
 Grothendieck, A. \emph{SGA 1--7}. Springer, 1971--1977.
\item
 Grothendieck, A. \emph{Tohoku Paper}. Tohoku Math. J., 1957.
\end{itemize}

\subsection{Textbooks and Surveys}\label{ch1966alexandergrothendieck:textbooks-and-surveys}

\begin{itemize}
\tightlist
\item
 Hartshorne, R. \emph{Algebraic Geometry}. Springer, 1977.
\item
 Milne, J. \emph{Étale Cohomology}. Princeton, 1980.
\item
 Eisenbud, D. and Harris, J. \emph{The Geometry of Schemes}. Springer,
 2000.
\item
 Vakil, R. \emph{The Rising Sea: Foundations of Algebraic Geometry}.
 Stanford, 2017.
\end{itemize}

\subsection{Biography and Exposition}\label{ch1966alexandergrothendieck:biography-and-exposition}

\begin{itemize}
\tightlist
\item
 Cartier, P. ``La folle journée de Grothendieck.'' \emph{Mathematical
 Intelligencer}, 1998.
\item
 Scharlau, W. \emph{Wer ist Alexander Grothendieck?} (biography).
\item
 Jackson, A. ``Alexand er Grothendieck.'' \emph{Notices of the AMS} 51,
 2004.
\end{itemize}

\begin{center}\rule{0.5\linewidth}{0.5pt}\end{center}

\emph{(End of report)}
